%%%
%%% Radical "⾢"
%%%
\section*{Radical 163: ``⾢'' (⻏)}\addcontentsline{toc}{section}{Radical 163: ⾢,⻏}\addcontentsline{loh}{figure}{\#\#\#\# 163: ⾢}

%%%%%%%%%% 邢 %%%%%%%%%%
\subsection*{邢}\addcontentsline{loh}{figure}{邢}

\begin{Entry}{邢}{6}{⾢}
  \begin{Phonetics}{邢}{xing2}
    \definition*{s.}{Usado em nomes de lugares, por exemplo XingTai (邢台) | Sobrenome: Xing}
  \seealsoref{邢台}{xing2tai2}
  \end{Phonetics}
\end{Entry}

\begin{Entry}{邢台}{6,5}{⾢,⼝}
  \begin{Phonetics}{邢台}{xing2tai2}
    \definition*{s.}{Xingtai, uma cidade na província de Hebei (河北)}
  \seealsoref{河北}{he2bei3}
  \end{Phonetics}
\end{Entry}

%%%%%%%%%% 那 %%%%%%%%%%
\subsection*{那}\addcontentsline{loh}{figure}{那}

\begin{Entry}{那}{6}{⾢}
  \begin{Phonetics}{那}{na1}
    \definition*{s.}{Sobrenome: Na}
  \end{Phonetics}
  \begin{Phonetics}{那}{na3}
    \definition{adv.}{expressa negação em perguntas retóricas}
    \definition{pron.}{qual? | qualquer que seja; qualquer que; para expressar incerteza em uma declaração | variante de 哪}
  \seealsoref{哪}{na3}
  \antonymref{这}{zhe4}
  \end{Phonetics}
  \begin{Phonetics}{那}{na4}[][HSK 1,2]
    \definition{conj.}{então; nessa situação; nesse caso; o mesmo que 那么}
    \definition{pron.}{aquele; aquilo; indica pessoas ou coisas distantes | aquele; aquilo; expressa muitas coisas, sem se referir especificamente a uma pessoa ou coisa, e é frequentemente usado em conjunto com 这}
  \seealsoref{那么}{na4me5}
  \seealsoref{这}{zhe4}
  \antonymref{这}{zhe4}
  \end{Phonetics}
  \begin{Phonetics}{那}{ne4}
    \definition{conj.}{então; nesse caso; o mesmo que 那么}
    \definition{pron.}{aquele; aquilo; pronúncia coloquial de 那 (\dpy{na4})}
  \seealsoref{那么}{na4me5}
  \antonymref{这}{zhe4}
  \end{Phonetics}
  \begin{Phonetics}{那}{nei4}
    \definition{conj.}{então; o mesmo que 那么}
    \definition{pron.}{aquele; aquilo; A pronúncia coloquial de 那 (\dpy{na4})}
  \seealsoref{那么}{na4me5}
  \antonymref{这}{zhe4}
  \end{Phonetics}
  \begin{Phonetics}{那}{nuo2}
    \definition*{s.}{Sobrenome: Nuo}
  \end{Phonetics}
\end{Entry}

\begin{Entry}{那儿}{6,2}{⾢,⼉}
  \begin{Phonetics}{那儿}{na4r5}[][HSK 1]
    \definition{pron.}{lá; ali; naquele lugar | então; naquela época (usado após 打, 从 e 由)}
  \seealsoref{从}{cong2}
  \seealsoref{打}{da3}
  \seealsoref{由}{you2}
  \synonymref{哪里}{na3li5}
  \synonymref{那边}{na4bian5}
  \synonymref{那里}{na4li3}
  \antonymref{这里}{zhe4li3}
  \antonymref{这儿}{zhe4r5}
  \end{Phonetics}
\end{Entry}

\begin{Entry}{那个}{6,3}{⾢,⼈}
  \begin{Phonetics}{那个}{na4ge5}
    \definition{pron.}{aquele | usado antes de verbos e adjetivos para indicar exagero | para substituir o discurso direto inconveniente}
  \end{Phonetics}
\end{Entry}

\begin{Entry}{那么}{6,3}{⾢,⼃}
  \begin{Phonetics}{那么}{na4me5}[][HSK 2]
    \definition{conj.}{então; nesse caso; afirmar o resultado esperado ou fazer um julgamento}
    \definition{pron.}{assim; dessa maneira; indica a natureza, o estado, a forma, o grau, etc. | assim; sobre; colocado antes do numeral, indica uma estimativa}
  \synonymref{那样}{na4yang4}
  \antonymref{这么}{zhe4me5}
  \end{Phonetics}
\end{Entry}

\begin{Entry}{那边}{6,5}{⾢,⾡}
  \begin{Phonetics}{那边}{na4bian5}[][HSK 1]
    \definition{pron.}{ali; acolá; aquele lado}
  \synonymref{那里}{na4li3}
  \synonymref{那儿}{na4r5}
  \antonymref{这边}{zhe4bian1}
  \antonymref{这里}{zhe4li3}
  \antonymref{这儿}{zhe4r5}
  \end{Phonetics}
\end{Entry}

\begin{Entry}{那会儿}{6,6,2}{⾢,⼈,⼉}
  \begin{Phonetics}{那会儿}{na4hui4r5}[][HSK 2]
    \definition{pron.}{então; naquela época; refere"-se ao passado ou ao futuro}
  \antonymref{这会儿}{zhe4hui4r5}
  \end{Phonetics}
\end{Entry}

\begin{Entry}{那时}{6,7}{⾢,⽇}
  \begin{Phonetics}{那时}{na4shi2}
    \definition{pron.}{então; naquela época; naqueles dias; geralmente se refere a um período de tempo distante do presente}
  \seealsoref{那时候}{na4 shi2hou5}
  \end{Phonetics}
\end{Entry}

\begin{Entry}{那时候}{6,7,10}{⾢,⽇,⼈}
  \begin{Phonetics}{那时候}{na4 shi2hou5}[][HSK 2]
    \definition{adv.}{naquela hora; em algum momento no passado}
  \seealsoref{那时}{na4shi2}
  \end{Phonetics}
\end{Entry}

\begin{Entry}{那里}{6,7}{⾢,⾥}
  \begin{Phonetics}{那里}{na4li3}[][HSK 1]
    \definition[行]{pron./s.}{lá; ali; aquele lugar; indica um lugar distante}
  \synonymref{何处}{he2chu4}
  \synonymref{哪里}{na3li5}
  \synonymref{哪儿}{na3r5}
  \synonymref{那边}{na4bian5}
  \synonymref{那儿}{na4r5}
  \antonymref{这边}{zhe4bian1}
  \antonymref{这里}{zhe4li3}
  \antonymref{这儿}{zhe4r5}
  \end{Phonetics}
\end{Entry}

\begin{Entry}{那些}{6,8}{⾢,⼆}
  \begin{Phonetics}{那些}{na4xie1}[][HSK 1]
    \definition{pron.}{aqueles; indica duas ou mais pessoas ou coisas}
  \synonymref{不计其数}{bu2ji4qi2shu4}
  \synonymref{大量}{da4liang4}
  \synonymref{大批}{da4pi1}
  \synonymref{众多}{zhong4duo1}
  \synonymref{诸多}{zhu1duo1}
  \end{Phonetics}
\end{Entry}

\begin{Entry}{那咱}{6,9}{⾢,⼝}
  \begin{Phonetics}{那咱}{na4 zan5}
    \definition{s.}{(informal) naquela época; então | (antigo) naquela época}
  \end{Phonetics}
\end{Entry}

\begin{Entry}{那样}{6,10}{⾢,⽊}
  \begin{Phonetics}{那样}{na4yang4}[][HSK 2]
    \definition{pron.}{assim; tal; desse tipo; desse gênero; dessa natureza; desse tipo; indica a natureza, o estado, a maneira, o grau ou refere"-se a uma ação ou situação específica}
  \synonymref{那么}{na4me5}
  \antonymref{这么}{zhe4me5}
  \antonymref{这样}{zhe4yang4}
  \end{Phonetics}
\end{Entry}

\begin{Entry}{那麽}{6,14}{⾢,⿇}
  \begin{Phonetics}{那麽}{na4me5}
    \variantof{那么}
  \end{Phonetics}
\end{Entry}

%%%%%%%%%% 邪 %%%%%%%%%%
\subsection*{邪}\addcontentsline{loh}{figure}{邪}

\begin{Entry}{邪}{6}{⾢}
  \begin{Phonetics}{邪}{xie2}[][HSK 7-9]
    \definition{adj.}{maligno; herético; estranho; irregular; impróprio | amaldiçoado; enfeitiçado; perverso; pessoas supersticiosas se referem às calamidades como sendo enviadas por fantasmas e deuses}
    \definition{s.}{fator patogênico; na medicina tradicional chinesa, as doenças são abordadas como causas de fatores ambientais}
  \antonymref{歪}{wai1}
  \antonymref{正}{zheng4}
  \end{Phonetics}
\end{Entry}

\begin{Entry}{邪恶}{6,10}{⾢,⼼}
  \begin{Phonetics}{邪恶}{xie2'e4}[][HSK 7-9]
    \definition{adj.}{mal; doente; perverso; vicioso}
    \definition{s.}{mal; doente; perverso; vicioso; refere"-se a pessoas ou forças malignas}
  \synonymref{凶恶}{xiong1'e4}
  \synonymref{凶狠}{xiong1hen3}
  \synonymref{罪恶}{zui4'e4}
  \antonymref{纯洁}{chun2jie2}
  \antonymref{善良}{shan4liang2}
  \antonymref{神圣}{shen2sheng4}
  \antonymref{正义}{zheng4yi4}
  \antonymref{正直}{zheng4zhi2}
  \end{Phonetics}
\end{Entry}

%%%%%%%%%% 邮 %%%%%%%%%%
\subsection*{邮}\addcontentsline{loh}{figure}{邮}

\begin{Entry}{邮}{7}{⾢}
  \begin{Phonetics}{邮}{you2}
    \definition*{s.}{Sobrenome: You}
    \definition{s.}{postal; correio; refere"-se a serviços postais | agência dos correios}
    \definition{v.}{postar; enviar pelo correio}
  \end{Phonetics}
\end{Entry}

\begin{Entry}{邮包}{7,5}{⾢,⼓}
  \begin{Phonetics}{邮包}{you2bao1}
    \definition[个]{s.}{encomenda postal; encomendas enviadas pelos correios}
  \synonymref{包裹}{bao1guo3}
  \synonymref{邮件}{you2jian4}
  \end{Phonetics}
\end{Entry}

\begin{Entry}{邮市}{7,5}{⾢,⼱}
  \begin{Phonetics}{邮市}{you2shi4}
    \definition{s.}{mercado de selos | mercado filatélico}
  \end{Phonetics}
\end{Entry}

\begin{Entry}{邮电}{7,5}{⾢,⽥}
  \begin{Phonetics}{邮电}{you2dian4}
    \definition{s.}{correios e telecomunicações; termo coletivo para serviços postais e de telecomunicações}
  \end{Phonetics}
\end{Entry}

\begin{Entry}{邮件}{7,6}{⾢,⼈}
  \begin{Phonetics}{邮件}{you2jian4}[][HSK 3]
    \definition[封,份,个,条]{s.}{correspondência; correio; assunto postal; termo que se refere a cartas, encomendas, etc., recebidos, transportados e entregues pelos correios | \emph{e"-mail}; refere"-se a \emph{e"-mails}, informações recebidas e enviadas através de caixas de correio eletrônico na \emph{Internet}, etc.}
  \end{Phonetics}
\end{Entry}

\begin{Entry}{邮局}{7,7}{⾢,⼫}
  \begin{Phonetics}{邮局}{you2ju2}[][HSK 4]
    \definition[家,个]{s.}{correio; agência dos correios; organizações que lidam com serviços postais}
  \end{Phonetics}
\end{Entry}

\begin{Entry}{邮政}{7,9}{⾢,⽁}
  \begin{Phonetics}{邮政}{you2zheng4}[][HSK 7-9]
    \definition{s.}{serviço postal; departamentos que gerenciam ou operam serviços de correio, telecomunicações, vale postal e distribuição de jornais/revistas}
  \synonymref{快递}{kuai4di4}
  \end{Phonetics}
\end{Entry}

\begin{Entry}{邮政编码}{7,9,12,8}{⾢,⽁,⽷,⽯}
  \begin{Phonetics}{邮政编码}{you2zheng4bian1ma3}
    \definition{s.}{código postal; CEP (Código de Endereçamento Postal)}
  \end{Phonetics}
\end{Entry}

\begin{Entry}{邮费}{7,9}{⾢,⾙}
  \begin{Phonetics}{邮费}{you2fei4}
    \definition{s.}{postagem | franquia postal}
    \definition{v.}{postar}
  \synonymref{邮资}{you2zi1}
  \end{Phonetics}
\end{Entry}

\begin{Entry}{邮迷}{7,9}{⾢,⾡}
  \begin{Phonetics}{邮迷}{you2mi2}
    \definition{s.}{colecionador de selos; filatelista}
  \end{Phonetics}
\end{Entry}

\begin{Entry}{邮资}{7,10}{⾢,⾙}
  \begin{Phonetics}{邮资}{you2zi1}
    \definition{s.}{postagem; a taxa cobrada pelos correios ao remetente da correspondência}
  \synonymref{邮费}{you2fei4}
  \synonymref{邮件}{you2jian4}
  \end{Phonetics}
\end{Entry}

\begin{Entry}{邮递}{7,10}{⾢,⾡}
  \begin{Phonetics}{邮递}{you2di4}
    \definition{s.}{entrega postal (por correio)}
    \definition{v.}{enviar por correio}
  \synonymref{发送}{fa1song5}
  \synonymref{快递}{kuai4di4}
  \synonymref{投递}{tou2di4}
  \antonymref{接收}{jie1shou1}
  \antonymref{领取}{ling3qu3}
  \end{Phonetics}
\end{Entry}

\begin{Entry}{邮票}{7,11}{⾢,⽰}
  \begin{Phonetics}{邮票}{you2piao4}[][HSK 3]
    \definition[枚,张,套,版]{s.}{selo; selo postal; comprovante vendido pelos correios, usado para colar nas correspondências para indicar que o porte foi pago}
  \end{Phonetics}
\end{Entry}

\begin{Entry}{邮编}{7,12}{⾢,⽷}
  \begin{Phonetics}{邮编}{you2bian1}[][HSK 7-9]
    \definition{s.}{código postal; código de área; abreviação de 邮政编码}
  \seealsoref{邮政编码}{you2zheng4bian1ma3}
  \end{Phonetics}
\end{Entry}

\begin{Entry}{邮箱}{7,15}{⾢,⾋}
  \begin{Phonetics}{邮箱}{you2xiang1}[][HSK 3]
    \definition{s.}{caixa de correio | \emph{mailbox}; refere"-se ao endereço de \emph{e"-mail}}
  \end{Phonetics}
\end{Entry}

%%%%%%%%%% 邻 %%%%%%%%%%
\subsection*{邻}\addcontentsline{loh}{figure}{邻}

\begin{Entry}{邻}{7}{⾢}
  \begin{Phonetics}{邻}{lin2}
    \definition{adj.}{vizinho; perto; adjacente; perto; próximo}
    \definition{s.}{vizinho | bairro; vizinhança}
  \end{Phonetics}
\end{Entry}

\begin{Entry}{邻国}{7,8}{⾢,⼞}
  \begin{Phonetics}{邻国}{lin2guo2}[][HSK 7-9]
    \definition{s.}{países vizinhos}
  \end{Phonetics}
\end{Entry}

\begin{Entry}{邻居}{7,8}{⾢,⼫}
  \begin{Phonetics}{邻居}{lin2ju1}[][HSK 5]
    \definition[个,位,名,家]{s.}{vizinho; pessoas ou famílias que moram muito perto}
  \end{Phonetics}
\end{Entry}

%%%%%%%%%% 郁 %%%%%%%%%%
\subsection*{郁}\addcontentsline{loh}{figure}{郁}

\begin{Entry}{郁}{8}{⾢}
  \begin{Phonetics}{郁}{yu4}
    \definition*{s.}{Sobrenome: Yu}
    \definition{adj.}{fortemente perfumado | luxuriante; exuberante | sombrio; deprimido}
  \end{Phonetics}
\end{Entry}

\begin{Entry}{郁郁葱葱}{8,8,12,12}{⾢,⾢,⾋,⾋}
  \begin{Phonetics}{郁郁葱葱}{yu4yu4cong1cong1}
    \definition{adj.}{exuberante e verde}
    \definition{expr.}{verdejante e exuberante; uma profusão selvagem de vegetação; luxuriantemente verde; ela cresce mais verde e mais fresca}
  \end{Phonetics}
\end{Entry}

%%%%%%%%%% 郊 %%%%%%%%%%
\subsection*{郊}\addcontentsline{loh}{figure}{郊}

\begin{Entry}{郊}{8}{⾢}
  \begin{Phonetics}{郊}{jiao1}
    \definition*{s.}{Sobrenome: Jiao}
    \definition{s.}{subúrbios; periferias; áreas ao redor da cidade}
  \end{Phonetics}
\end{Entry}

\begin{Entry}{郊区}{8,4}{⾢,⼖}
  \begin{Phonetics}{郊区}{jiao1qu1}[][HSK 5]
    \definition[个,片,块]{s.}{subúrbios; arredores; periferia; área ao redor da cidade que está administrativamente sob a jurisdição da cidade}
  \end{Phonetics}
\end{Entry}

\begin{Entry}{郊外}{8,5}{⾢,⼣}
  \begin{Phonetics}{郊外}{jiao1wai4}[][HSK 7-9]
    \definition{s.}{subúrbio; periferia; a zona rural ao redor de uma cidade; a área fora da cidade (referindo"-se a uma cidade específica)}
  \end{Phonetics}
\end{Entry}

\begin{Entry}{郊游}{8,12}{⾢,⽔}
  \begin{Phonetics}{郊游}{jiao1you2}[][HSK 7-9]
    \definition{s.}{passeio; excursão}
  \end{Phonetics}
\end{Entry}

%%%%%%%%%% 郑 %%%%%%%%%%
\subsection*{郑}\addcontentsline{loh}{figure}{郑}

\begin{Entry}{郑}{8}{⾢}
  \begin{Phonetics}{郑}{zheng4}
    \definition*{s.}{O estado de Zheng durante o período dos Reinos Combatentes; abreviatura de 郑州 | Sobrenome: Zheng}
  \seealsoref{郑州}{zheng4zhou1}
  \end{Phonetics}
\end{Entry}

\begin{Entry}{郑州}{8,6}{⾢,⼮}
  \begin{Phonetics}{郑州}{zheng4zhou1}
    \definition*{s.}{Zhengzhou, capital da província de Henan (河南)}
  \seealsoref{河南}{he2nan2}
  \end{Phonetics}
\end{Entry}

\begin{Entry}{郑重}{8,9}{⾢,⾥}
  \begin{Phonetics}{郑重}{zheng4zhong4}[][HSK 7-9]
    \definition{adj.}{sério; solene; fervoroso}
  \synonymref{谨慎}{jin3shen4}
  \synonymref{认真}{ren4zhen1}
  \synonymref{慎重}{shen4zhong4}
  \synonymref{严肃}{yan2su4}
  \end{Phonetics}
\end{Entry}

%%%%%%%%%% 部 %%%%%%%%%%
\subsection*{部}\addcontentsline{loh}{figure}{部}

\begin{Entry}{部}{10}{⾢}
  \begin{Phonetics}{部}{bu4}[][HSK 3]
    \definition*{s.}{Sobrenome: Bu}
    \definition{clas.}{usado para obras de literatura, livros, filmes, etc.}
    \definition[根]{s.}{parte; seção | unidade; ministério; departamento; conselho | sede; matriz; quartel general | tropas; forças | divisão; região}
    \definition{v.}{comandar; liderar}
  \end{Phonetics}
\end{Entry}

\begin{Entry}{部下}{10,3}{⾢,⼀}
  \begin{Phonetics}{部下}{bu4xia4}
    \definition{s.}{tropas sob o comando de alguém; pessoas sob o comando do exército | subordinado}
  \synonymref{部属}{bu4shu3}
  \synonymref{下级}{xia4ji2}
  \synonymref{下属}{xia4shu3}
  \antonymref{上级}{shang4ji2}
  \antonymref{上司}{shang4si5}
  \end{Phonetics}
\end{Entry}

\begin{Entry}{部门}{10,3}{⾢,⾨}
  \begin{Phonetics}{部门}{bu4men2}[][HSK 3]
    \definition[个]{s.}{departamento; ramo; classe; seção; partes ou unidades que compõem um todo}
  \end{Phonetics}
\end{Entry}

\begin{Entry}{部分}{10,4}{⾢,⼑}
  \begin{Phonetics}{部分}{bu4fen5}[][HSK 2]
    \definition[个,些,快,份]{s.}{parte; seção; porção; parte do todo; alguns indivíduos dentro do todo | ramo; parte separada de um sistema ou entidade}
  \seealsoref{局部}{ju2bu4}
  \synonymref{个别}{ge4bie2}
  \synonymref{局部}{ju2bu4}
  \antonymref{全部}{quan2bu4}
  \antonymref{全体}{quan2ti3}
  \antonymref{一切}{yi2qie4}
  \antonymref{整个}{zheng3ge4}
  \antonymref{整体}{zheng3ti3}
  \antonymref{总体}{zong3ti3}
  \end{Phonetics}
\end{Entry}

\begin{Entry}{部长}{10,4}{⾢,⾧}
  \begin{Phonetics}{部长}{bu4zhang3}[][HSK 3]
    \definition[个,位,名]{s.}{ministro; chefe de departamento; um alto funcionário do estado encarregado pelo chefe de estado ou chefe executivo do governo da gestão das atividades governamentais de um departamento | chefe de seção; líder tribal}
  \end{Phonetics}
\end{Entry}

\begin{Entry}{部队}{10,4}{⾢,⾩}
  \begin{Phonetics}{部队}{bu4dui4}[][HSK 6]
    \definition[支,个]{s.}{militar; exército; forças armadas | tropas; refere"-se a uma parte do exército}
  \end{Phonetics}
\end{Entry}

\begin{Entry}{部件}{10,6}{⾢,⼈}
  \begin{Phonetics}{部件}{bu4jian4}[][HSK 7-9]
    \definition[个]{s.}{peças; partes; componentes; um componente de uma máquina, montado a partir de várias partes | partes; componentes (para caracteres chineses); uma unidade de caracteres chineses composta por traços, por exemplo, 氵, 礻, 口 são todos componentes de caracteres chineses}
  \end{Phonetics}
\end{Entry}

\begin{Entry}{部位}{10,7}{⾢,⼈}
  \begin{Phonetics}{部位}{bu4wei4}[][HSK 5]
    \definition{s.}{lugar; posição (usado principalmente para o corpo humano)}
  \end{Phonetics}
\end{Entry}

\begin{Entry}{部族}{10,11}{⾢,⽅}
  \begin{Phonetics}{部族}{bu4zu2}
    \definition{adj.}{tribal}
    \definition{s.}{tribo; clã; grupo de famílias}
  \synonymref{民族}{min2zu2}
  \end{Phonetics}
\end{Entry}

\begin{Entry}{部属}{10,12}{⾢,⼫}
  \begin{Phonetics}{部属}{bu4shu3}
    \definition{adj.}{subordinado | afiliado a um ministério}
    \definition{s.}{tropas sob o comando de alguém | unidades subordinadas departamentais}
  \synonymref{部下}{bu4xia4}
  \synonymref{下属}{xia4shu3}
  \antonymref{领袖}{ling3xiu4}
  \antonymref{上级}{shang4ji2}
  \antonymref{上司}{shang4si5}
  \end{Phonetics}
\end{Entry}

\begin{Entry}{部署}{10,13}{⾢,⽹}
  \begin{Phonetics}{部署}{bu4shu3}[][HSK 7-9]
    \definition{v.}{organizar; implantar; dispor; organizar ou dispor de maneira planejada (usado principalmente em grandes aspectos)}
  \end{Phonetics}
\end{Entry}

%%%%%%%%%% 都 %%%%%%%%%%
\subsection*{都}\addcontentsline{loh}{figure}{都}

\begin{Entry}{都}{10}{⾢}
  \begin{Phonetics}{都}{dou1}[][HSK 1]
    \definition{adv.}{todos; representa a soma total | apenas por causa de; usado em conjunto com a palavra 是, explica o motivo | mesmo; até; indicativo de ênfase | já; significa 已经}
  \seealsoref{皆}{jie1}
  \seealsoref{是}{shi4}
  \seealsoref{也}{ye3}
  \seealsoref{已经}{yi3jing1}
  \synonymref{皆}{jie1}
  \synonymref{全}{quan2}
  \synonymref{也}{ye3}
  \end{Phonetics}
  \begin{Phonetics}{都}{du1}
    \definition*{s.}{Sobrenome: Du}
    \definition[座]{s.}{capital | cidade grande; metrópole}
  \end{Phonetics}
\end{Entry}

\begin{Entry}{都市}{10,5}{⾢,⼱}
  \begin{Phonetics}{都市}{du1shi4}[][HSK 6]
    \definition[个]{s.}{cidade grande; grandes cidades}
  \end{Phonetics}
\end{Entry}

\begin{Entry}{都会}{10,6}{⾢,⼈}
  \begin{Phonetics}{都会}{du1hui4}[][HSK 7-9]
    \definition{s.}{cidade; metrópole}
  \end{Phonetics}
\end{Entry}

%%%%%%%%%% 鄙 %%%%%%%%%%
\subsection*{鄙}\addcontentsline{loh}{figure}{鄙}

\begin{Entry}{鄙}{13}{⾢}
  \begin{Phonetics}{鄙}{bi3}
    \definition*{s.}{Sobrenome: Bi}
    \definition{adj.}{baixo; mesquinho; vulgar | rústico; básico; desprezível}
    \definition{pron.}{(auto"-depreciativo) meu}
    \definition{s.}{Literário: um lugar remoto; cidade fronteiriça; cidade pequena}
    \definition{v.}{Literário: desprezar; desdenhar; menosprezar; olhar de cima para baixo}
  \end{Phonetics}
\end{Entry}

\begin{Entry}{鄙视}{13,8}{⾢,⾒}
  \begin{Phonetics}{鄙视}{bi3shi4}[][HSK 7-9]
    \definition{v.}{desprezar; desdenhar; menosprezar}
  \synonymref{忽视}{hu1shi4}
  \synonymref{看不起}{kan4bu5qi3}
  \synonymref{歧视}{qi2shi4}
  \synonymref{轻蔑}{qing1mie4}
  \synonymref{小看}{xiao3kan4}
  \antonymref{崇拜}{chong2bai4}
  \antonymref{敬佩}{jing4pei4}
  \antonymref{看好}{kan4hao3}
  \antonymref{看重}{kan4zhong4}
  \antonymref{羡慕}{xian4mu4}
  \antonymref{欣赏}{xin1shang3}
  \antonymref{重视}{zhong4shi4}
  \antonymref{尊敬}{zun1jing4}
  \antonymref{尊重}{zun1zhong4}
  \end{Phonetics}
\end{Entry}

%%%%% EOF %%%%%

